%! Modeling with Exponential Functions: Interest, Half-Life, and e — Unit 7, Lesson 7.4 — Warm-Up
\documentclass[10pt]{article}
\usepackage{algebra2-article}
\usepackage{algebra2-boxes}
\newcommand{\TallMath}[1]{$\displaystyle #1\rule[-1.4em]{0pt}{3.2em}$}

\begin{document}

\pageheader{Unit 7, Lesson 7.4}{Warm-Up}
\namedateperiod

\vspace{0.10in}

\begin{notesbox}{Getting Ready: Skills We'll Use Today}
\small Today's exponential models come from money and medicine. Nothing below is new: item 1 is Lesson
7.1's percent-to-factor translation, item 2 is Unit 6 exponent work, and item 3 is yesterday.

\vspace{5pt}
\textbf{1. One yearly rate, cut into smaller pieces.} An account pays $6\%$ per year. If interest is
added more than once a year, each payment uses only \emph{part} of that rate.

\vspace{3pt}
\begin{center}
\renewcommand{\arraystretch}{1.4}
\begin{tabular}{|l|c|c|c|c|}
\hline
\textbf{Interest is added} & \textbf{times per year} & \textbf{rate per period} & \textbf{factor per period} & \textbf{periods in 5 years} \\
\hline
once a year & $1$ & $0.06$ & $1.06$ & $5$ \\
\hline
every quarter & $4$ & \blank{1.5cm} & \blank{1.5cm} & \blank{1.5cm} \\
\hline
every month & $12$ & \blank{1.5cm} & \blank{1.5cm} & \blank{1.5cm} \\
\hline
\end{tabular}
\end{center}

\vspace{0.20in}

\textbf{2. Halving, over and over.} Evaluate. Round decimals to three places.

\vspace{4pt}
\hspace*{1.2em} $\left(\dfrac12\right)^{3} =$ \blank{1.6cm} \qquad
$\left(\dfrac12\right)^{0} =$ \blank{1.6cm} \qquad
$\left(\dfrac12\right)^{1/2} \approx$ \blank{1.6cm} \qquad
$\left(\dfrac12\right)^{1.5} \approx$ \blank{1.6cm}

\vspace{6pt}
\hspace*{1.2em} A medicine loses half of whatever is left every $6$ hours. How many \emph{halvings}
happen in $18$ hours? \blank{1.6cm}

\vspace{4pt}
\hspace*{1.2em} How many halvings happen in $9$ hours? \blank{1.6cm} \; Explain how a number of
halvings can fail to be a whole number. \writeline

\vspace{0.22in}

\textbf{3. From yesterday.} Solve using a common base (Lesson 7.3).

\vspace{4pt}
\hspace*{1.2em} $2^{\,t} = 32$ \hfill $t=$ \blank{2.0cm}

\vspace{6pt}
\hspace*{1.2em} Now look at $\;(1.05)^{\,t} = 2$. \; What would Step 1 ask you to do, and why can you
not do it? \writeline

\end{notesbox}

\end{document}
